 %%%%%%%%%%%%%%%%%%%%%%%%%%%%%%%%%%%%%%%%%%%%%%%%%%%%%%%%%%%%%%%%%%%%%%%%%%%%%%%%%%%%%%
% Template fuer Abschlussarbeiten von Studierenden der  Frankfurt University of Applied Sciences
% 
% erstellt von: Prof. Dr.-Ing. Thomas Hollstein
% 
% Last revision: 22.06.2023
%
%%%%%%%%%%%%%%%%%%%%%%%%%%%%%%%%%%%%%%%%%%%%%%%%%%%%%%%%%%%%%%%%%%%%%%%%%%%%%%%%%%%%%%

% Grundlegende Formatierungen
\documentclass[
	%twoside, 
	%openright,
	titlepage,
	numbers=noenddot,
	headinclude,
	%1headlines,
	footinclude=true,
	%cleardoublepage=empty,
	%BCOR=5mm,
	fontsize=12pt,%11pt,
	paper=a4paper,
	%letterpaper
	ngerman,
	american,
	%table
 ]
% Dokumententyp: (legt grundlegende Formatierrichtlinien fest)
%{report}
{scrreprt}
%{book}

% Farbigen Text
\usepackage{xcolor}
\usepackage{colortbl}

\parindent = 10pt  % Neue Absaetze nicht einrueckenn
\parskip = 1ex % Neue Absaetze: 1/2 Zeile Abstand


%%%%%%%%%%%%%%%%%%%%%%%%%%%%%%%%%%%%%%%%%%%%%%%%%%%%%%%%%%%%%%%%%%%%%%%%%%%%%%%%%%%%%%


%%%%%%%%%% zu nutzende Pakete einbinden: %%%%%%%%%%

%%%%% Package Kommentare:
\usepackage{comment}
\usepackage[colorinlistoftodos]{todonotes}
\newcommand{\mycomment}[1]{\todo[inline,linecolor=green,backgroundcolor=yellow!25,bordercolor=green,caption={}]{todo: #1}}

%%%%% package Sprachunterstuetzung:
\usepackage[utf8]{inputenc} % Eingabekodierung auf UTF-8 setzen
\usepackage[T1]{fontenc}    % Schriftkodierung auf T1 setzen
\usepackage[ngerman]{babel}
\usepackage{csquotes}

%%%%% Package Multirow ermoeglicht, dass sich ein Kaestchen in einer Tabelle ueber mehrere Zeilen erstrecken kann:
\usepackage{multirow}
\usepackage{array}
\usepackage{ragged2e}
\usepackage{booktabs}
\usepackage{arydshln}
% Beispiele unter: https://texblog.org/2012/12/21/multi-column-and-multi-row-cells-in-latex-tables/
%Vereinigung von Feldern in Tabellen:
%\multicolumn{number cols}{align}{text} % align: l,c,r
%\multirow{number rows}{width}{text}

%%%% Serifenloser Textstil fuer das ganze Dokument: 
\RequirePackage[sfdefault,lf]{carlito}
\usepackage{lmodern}
% \RequirePackage[T1]{fontenc}
% to imitate Calibri:
\renewcommand*\familydefault{\sfdefault} %% Base font of the document is to be sans serif
\usepackage{lipsum}

%%%%%% Grafiken einbinden
\usepackage{float}
\usepackage{graphicx}
\usepackage{subcaption}
\usepackage{lscape}
\usepackage{svg}
\usepackage{pgf-pie}
\usepackage{transparent}
\usepackage{pdfpages}
\usepackage{pgfplots}
\pgfplotsset{compat=newest}
\pgfplotsset{plot coordinates/math parser=false}
\usetikzlibrary{patterns,arrows,shapes,positioning,shadows,trees,calc,positioning,external}
\usepackage{tikz}
\tikzexternalize[prefix=Figures/tikzexternalize_cache/,optimize command away=\includepdf]
%\tikzexternaldisable
% https://golatex.de//wiki/%5cincludegraphics

%%%%%% Fuer Baumdiagramme
\tikzset{
	basic/.style  = {draw, rectangle},
	root/.style   = {basic, rounded corners=2pt, thin, align=center},
	level 2/.style = {basic, rounded corners=6pt, thin,align=center, text width=5em},
	level 3/.style = {basic, align=left, draw=none, text width=60}
}

%%%%%% Mathematik-Paket AMSMath:
\usepackage[fleqn,reqno]{amsmath}
\usepackage{mathtools}
%%%%%% Brueche
\usepackage{nicefrac} % For comparison
\usepackage{xfrac}    % Works better with other fonts
\usepackage{romannum}

%%%%%%%%%% Seitengeometrie einstellen: %%%%%%%%%%
% https://tex.stackexchange.com/questions/344241/logo-as-header-using-fancyhdr-package
\usepackage{geometry}
\geometry{verbose,
	bmargin=2cm,
	lmargin=2.8cm,
	rmargin=3.5cm
	%footskip=-25pt
}

%%%%%% die Höhe des Top-Margin errechnen:
\newlength\mytopmargin
\newsavebox{\headbox}\savebox{\headbox}{
	%\includegraphics[width=0.3\textwidth]{Figures/xxx.png} \hfill
	\raisebox{-1ex}{\includegraphics[width=0.1\textwidth]{Figures/pdf/FraUas_Logo.pdf}}
}    
\setlength{\mytopmargin}{\totalheightof{\usebox{\headbox}}+2cm}
\geometry{verbose,
	tmargin=\mytopmargin,
	headheight=1.1\mytopmargin,
	footskip=6ex
}

% Flexiblere Tabellen
\usepackage{tabularx}
% \usepackage{xltabular}
\def\tabularxcolumn#1{m{#1}}
\usepackage{enumitem}
%\usepackage{enumerate}
\usepackage{layouts}
\usepackage{listliketab}
\usepackage{rotating} 

% Boxen für Lickert Skalen
\usepackage{wasysym}
\newcommand\insq[1]{%
	\Square\ #1\quad%
}

% Durchstreichen von Text
\usepackage[normalem]{ulem} % mit \sout

%%%%%% Gesamtseitenzahl verwenden:
\usepackage{totpages}

%%%%%% Kalkulationen:
%https://tex.stackexchange.com/questions/30081/how-can-i-sum-two-values-and-store-the-result-in-other-variable
\usepackage{tikz}
\usetikzlibrary{math}
\usepackage{amsmath}
\usepackage{amssymb}

%%%%%% Gestaltung von Kopf- und Fusszeilen:
% https://tex.stackexchange.com/questions/344241/logo-as-header-using-fancyhdr-package
\usepackage{fancyhdr}
\pagestyle{fancy}  % Eigener Seitenstil
\fancyhf{}         % Alle Kopf- und Fußzeilenfelder bereinigen
%\fancyhead[L]{} % Kopfzeile links
\fancyhead[l]{\leftmark
	%\makebox[0.3\textwidth]{\includegraphics[width=0.3\textwidth]{Figures/xxx.png}}
}  
%\fancyhead[c]{\hspace*{0.15\textwidth}\rightmark}
%\fancyhead[C]{\usebox\headbox}                        % Zentrierte Kopfzeile
\fancyhead[R]{
	\makebox[0.2\textwidth]{\raisebox{-5ex}{\hspace*{9.5ex}\includegraphics[width=0.1\textwidth]{Figures/pdf/FraUas_Logo.pdf}}}
}  % Kopfzeile rechts
\renewcommand{\headrulewidth}{0.4pt} % Obere Trennlinie
%\fancyfoot[L]{\today}
\fancyfoot[C]{\ThesisTitleShort} 
%\fancyfoot[R]{Seite \thepage ~von \ref{TotPages}}  % Seitennummer
\fancyfoot[R]{\thepage} % Seitennummer 
\renewcommand{\footrulewidth}{0.4pt} % Untere Trennlinie
\setlength{\mytopmargin}{\totalheightof{\usebox\headbox} +2cm}
%Unterschied zwischen geraden/ungeraden Seiten:
%\fancyhead[OR]{} % "O" steht für "odd", also ungerade Seiten
%\fancyhead[ER]{} % "E" für "even", also gerade Seiten.

%%%%% Erweiterte Formate für Listen/Aufzaehlungen:
\usepackage{paralist}
%Default-Items fuer die vier moeglichen Verschachtelungsebenen:
\setdefaultitem{}{\textbullet}{$\star$}{}

%%%%% FRA-UAS CI Farben:
%\definecolor{airforceblue}{rgb}{0.36, 0.54, 0.66}
% CI-Farben FRA-UAS (blau):
\definecolor{FRAUAS_Blue_Dark}{RGB}{45, 137, 204}
\definecolor{FRAUAS_Blue_Light}{RGB}{182, 210, 228}
% CI-Farben FRA-UAS: FB1
\definecolor{FRAUAS_FB1_Dark}{RGB}{124, 128, 52}
\definecolor{FRAUAS_FB1_Light}{RGB}{213, 213, 179}
% CI-Farben FRA-UAS: FB2
\definecolor{FRAUAS_FB2_Dark}{RGB}{255, 158, 27}
\definecolor{FRAUAS_FB2_Light}{RGB}{251, 221, 173}
% CI-Farben FRA-UAS: FB3
\definecolor{FRAUAS_FB3_Dark}{RGB}{196, 213, 42}
\definecolor{FRAUAS_FB3_Light}{RGB}{237, 240, 166}
% CI-Farben FRA-UAS: FB4
\definecolor{FRAUAS_FB4_Dark}{RGB}{204, 31, 47}
\definecolor{FRAUAS_FB4_Light}{RGB}{240, 166, 183}

%%%%% Sektionstitel nach CI-Farben einfärben:
\usepackage{titlesec}
\titleformat{\section}
{\color{blue50}\normalfont\Large\bfseries} %Titel
{\color{blue50}\thesection}{1em}{}
%{\color{FRAUAS_Blue_Dark}\normalfont\Large\bfseries} %Titel
%{\color{FRAUAS_Blue_Dark}\thesection}{1em}{}
\titleformat{\subsection}
{\color{blue50}\normalfont\large\bfseries} %Titel
{\color{blue50}\thesubsection}{1em}{}
%{\color{FRAUAS_Blue_Dark}\normalfont\large\bfseries} %Titel
%{\color{FRAUAS_Blue_Dark}\thesubsection}{1em}{}
\titleformat{\subsubsection}
{\color{black}\normalfont\bfseries} %Titel
{\color{black}\thesubsection}{0.5em}{}

%\usepackage{appendix}

%%%%% Erweiterte Bibliographie-Stile:
%\usepackage{harvard}
% Title Page
% (wir generieren den Titel per Handlayout und verwenden
% daher die folgenden Befehle nicht)
%\title{}
%\author{}

% https://www.overleaf.com/learn/latex/Hyperlinks
\usepackage{hyperref}
\hypersetup{
	%hyperindex=true,
	%linktocpage=true, % Seitenzahl statt Titel verlinkt
	colorlinks=true,
	linkcolor=black,
	filecolor=black,      
	urlcolor=black,
	pdftitle={Thesis},
	citecolor=black,
	pdfpagemode=UseNone, % Keine automatische Ansichtseinstellung
	pdfstartpage=1,      % Startet auf der ersten Seite
	pdfstartview=Fit     % Zeigt die gesamte Seite im Fenster
}

%\usepackage[hyphens]{url}  %% stellt \url{} zur Verfuegung

%%%%% modernes BibLaTeX mit biber %%%%%
% https://golatex.de/viewtopic.php?t=13917
%\usepackage[style=ieee-alphabetic,backend=biber, natbib=true]{biblatex}

% Literatur nach Erscheinen im Text sortiert
%\usepackage[sorting=none, style=numeric,backend=biber, natbib=true]{biblatex}
%\usepackage[style=apa, backend=biber, natbib=true, sorting=nyt, sortcites=false]{biblatex}
\usepackage[style=ieee, backend=biber, sorting=none, natbib=true]{biblatex}
%\usepackage[style=numeric, backend=biber, natbib=true, sorting=nyt, sortcites=false]{biblatex}
% Ohne eingestellte Sortierung
%\usepackage[ style=numeric,
%backend=biber, natbib=true]{biblatex}

\addbibresource{bibliography.bib}

% Immer shortautor anzeigen, falls vorhanden
\makeatletter
\def\cbx@apa@ifnamesaved{\@firstoftwo}
\makeatother

% Zeilenabstand für das ganze Dokument (1.0 = Normalwert):
\renewcommand{\baselinestretch}{1.0}

%%%%% Acronymdefinitionen einbinden:
\usepackage[printonlyused]{acronym}
\usepackage{tocloft}
%\input{TeXFiles/005_Acronyms}

%%%%% Glossar einbinden
\usepackage{glossaries}
\makeglossaries

%% Glossareintraege
\newglossaryentry{Pruefung}{
	name=Prüfung,
	description={Eine Prüfung enthält mehrere Einzelmessungen und kann entweder Batch-Sequenziell oder DUT-Sequenziell durchgeführt werden}
}
%\usepackage[record]{bib2gls}
%\addbibresource{Glossar.bib}

\setcounter{tocdepth}{1} % Inhaltsverzeichnis bis zur 2. Nummerierungstiefe


%%%%%%%%%%%%%%%%%%%%%%%%%%%%%%%%%%%%%%%%%%%%%%%%%%%%%%%%%%%%%%%%%%%%%%%%%%%%%%%%%%%%%%


%%%%%%%%%%%%%%%%%%%%%%%%%%%%%%%%%%%%%%%%%%%%
%%%%% hier beginnt das eigentliche Dokument 
%%%%%%%%%%%%%%%%%%%%%%%%%%%%%%%%%%%%%%%%%%%%

\begin{document}
	%%%%% Settings einbinden:
	\newcommand{\myName}{Paul Glaser}
\newcommand{\ThesisTitle}{Zuverlässigkeitsuntersuchung elektrischer Traktionsmotoren}
\newcommand{\ThesisTitleShort}{Untersuchung el. Traktionsmotoren}
%\newcommand{\ThesisSubtitle}{This is the subtitle of the thesis}
\newcommand{\ThesisDegree}{Bachelor of Engineering (B. Eng.)} 
%\newcommand{\ThesisDegree}{Master of Science (M. Sc.)}
\newcommand{\myStudentId}{1346190}
\newcommand{\Supervisor}{Prof. Dr.-Ing. Thomas Hollstein}
\newcommand{\CoSupervisor}{??}
\newcommand{\CompSupervisor}{M.Sc. Timo Petri}
\newcommand{\Faculty}{Faculty of Computer Science and Engineering}
\newcommand{\University}{Frankfurt University of Applied Sciences}
\newcommand{\UniversityLocation}{Frankfurt}
\newcommand{\ThesisDeliveryDate}{14. Dezember 2024 ??}
\newcommand{\CompanyName}{Robert Bosch GmbH}  % don't modify this line, if no company is involved

\newcommand{\VarPicWidthA}{\textwidth}
\newcommand{\VarPicWidthB}{0.9\textwidth}
\newcommand{\VarPicWidthC}{0.8\textwidth}
\newcommand{\VarPicWidthD}{0.7\textwidth}
\newcommand{\VarPicWidthE}{0.6\textwidth}
\newcommand{\VarPicWidthF}{0.49\textwidth}
\newcommand{\VarPicWidthG}{0.32\textwidth}
\newcommand{\VarPicWidthH}{0.24\textwidth}

\definecolor{blue40}{rgb}{0, 0.3843, 0.6039}
\definecolor{blue55}{rgb}{0, 0.5333, 0.8314}
\definecolor{green40}{rgb}{0, 0.4235, 0.2275}
\definecolor{green55}{rgb}{0.1294, 0.5843, 0.3412}
\definecolor{red40}{rgb}{0.7451, 0, 0.0157}
\definecolor{red55}{rgb}{1.0000, 0.1294, 0.1412}

\newcommand{\ti}[1]{_\mathrm{#1}}
\newcommand\tab[1][1cm]{\hspace*{#1}}

% Definieren von Stilen für die Formen und Pfeile im Diagramm
\tikzstyle{startstop} = [rectangle, rounded corners, minimum width=3cm, minimum height=1cm,text centered, draw=black]
\tikzstyle{process} = [rectangle, text centered, draw=black]
\tikzstyle{info} = [rectangle, text centered, draw=none]
\tikzstyle{decision} = [diamond, minimum width=3cm, minimum height=1cm, text centered, draw=black]
\tikzstyle{arrow} = [thick,->,>=stealth]
\tikzstyle{vorangeganen} = [rectangle, text centered, draw=black, fill=blue55, fill opacity=0.1]
\tikzstyle{hier} = [rectangle, text centered, draw=black]
\tikzstyle{weiterführen} = [rectangle, text centered, draw=black, fill=green55, fill opacity=0.1]


%%%%%%%%%%%%%%%%%%%%% FOR ENGLISH LANGUAGE THESIS: %%%%%%%%%%%%%%%%%%%%

% By default the thesis language is german
% if you want to set it to ENGLISH, then UNCOMMENT the FOLLOWING LINE by removing the leading "%":

%\newcommand*{\ThesisLanguageIsEnglish}{}

%% End of File
	\sloppy %Formatierungsueberstaende am Zeilenende vermeiden
	% https://latexref.xyz/_005cfussy-_0026-_005csloppy.html

	\frenchspacing  
	%Ein Leerzeichen nach Satzende
	%https://texwelt.de/fragen/1154/was-ist-french-spacing-was-macht-frenchspacing

	%\raggedbottom   
	%Standard is \flushbottom, dass heisst alle Seiten werden so gedehnt, dass sie  gleich hoch sind 
	%Schaltet man \raggedbottom ein, ist dies nicht so

	\ifdefined\ThesisLanguageIsEnglish
		\selectlanguage{american}
	\else
		\selectlanguage{ngerman} % ngerman, american
	\fi
	%Deutsch nach neuer Rechtschreibung als Standardsprache fuer das Dokument einstellen

	%\renewcommand*{\bibname}{new name}
	%\setbibpreamble{}

	% Numerierungstiefe setzen:
	% 2: bis subsection (Standard)
	% 3: bis subsubsection
	\setcounter{secnumdepth}{2} %setzt die Numerierungstiefe



	%\renewcommand{\thepage}{\Roman{page}}
	\pagestyle{plain} % Seite ohne Kopf- und Fusszeilen darstellen



	% hier wird die Titelseite eingebunden:
	%*******************************************************
% Titlepage
%*******************************************************
%%%
%%% title page (german)
%%%
\thispagestyle{empty}
\pdfbookmark[0]{Titelblatt}{title}

\begin{titlepage}
	
	% Fra-UAS Logo
	\vspace*{-2,5cm}
  	\begin{center}
    		\includegraphics[width=7.7cm]{Figures/pdf/FraUas_Logo.pdf} \\ 
  	\end{center}

	% Name FH und Fakultaet
  	\begin{center}
		\vspace{0.1cm}
		\LARGE \textbf{\University}\\
		\vspace{0.4cm}
		\Large -- \Faculty --
	\end{center}
	
	\vfill
	
	% Titel (--> Bitte in TeXFiles/000_Settings eintragen)
	\begin{center}
		\huge \textbf{\ThesisTitle}\\
		\vspace{0.4cm}
		\LARGE \ThesisSubtitle
	\end{center}
	
	\vfill
	
	% Untertitel (--> Bitte in TeXFiles/000_Settings eintragen)
	\ifdefined\ThesisLanguageIsEnglish
		\begin{center}
			\Large Thesis submitted in order to obtain the academic degree\\
			\vspace{0.3cm}
			\Large \ThesisDegree
		\end{center}
	\else
		\begin{center}
			\Large Abschlussarbeit zur Erlangung des akademischen Grades\\
			\vspace{0.3cm}
			\Large \ThesisDegree   %%%%% >>>>> Bitte in TeXFiles/000_Settings eintragen
		\end{center}
	\fi
	
	\vfill
	
	% Datum, Name, Matrikelnummer (--> Bitte in TeXFiles/000_Settings eintragen)
	\ifdefined\ThesisLanguageIsEnglish 
  		\begin{center}
    			\Large submitted on \ThesisDeliveryDate\ on\\
    			\vspace{0.3cm}
    			\Large \textbf{\myName}\\
    			\vspace{0.3cm}
    			\normalsize Student ID: \myStudentId
  		\end{center}
  	\else
		\begin{center}
			\Large vorgelegt am \ThesisDeliveryDate\ von\\
			\vspace{0.3cm}
			\Large \textbf{\myName}\\
			\vspace{0.3cm}
			\normalsize Matrikelnummer: \myStudentId
		\end{center}
  	\fi
	
	\vfill
	
	% Betreuer (--> Bitte in TeXFiles/000_Settings eintragen)
	\ifdefined\ThesisLanguageIsEnglish 
		\begin{center}
			\begin{tabular}{lll}
				First Supervisor    & : & \Supervisor \\     %%%%% >>>>> Bitte in TeXFiles/000_Settings eintragen
				Second Supervisor & : & \CoSupervisor\\    %%%%% >>>>> Bitte in TeXFiles/000_Settings eintragen
			\end{tabular}
		\end{center}
	\else
		\begin{center}
			\begin{tabular}{lll}
				Erstprüfer    		& : & \Supervisor \\     %%%%% >>>>> Bitte in TeXFiles/000_Settings eintragen
				Zweitprüfer			& : & \CoSupervisor\\    %%%%% >>>>> Bitte in TeXFiles/000_Settings eintragen
				Betreuer Fa. Bosch 	& : & \CompSupervisor\\    %%%%% >>>>> Bitte in TeXFiles/000_Settings eintragen
			\end{tabular}
		\end{center} 
	\fi
	
	\newpage
	
\end{titlepage}
 
 %% End of File
	%*******************************************************
% Non Disclosure Notice
%*******************************************************

\ifdefstring{\CompanyName}{NoCompanyInvolved}{}{
	\ifdefined\ThesisLanguageIsEnglish
		\chapter*{Non-Disclosure Notice}
			\thispagestyle{empty}
			This scientific thesis contains information of the company \CompanyName. The disclosure or
			unauthorized use of this thesis as a whole or in parts is prohibited. The prohibition shall only
			apply to information pertaining to the company \CompanyName. Any other exceptions require the
			written consent of \CompanyName.
		\medskip
	\else
		\chapter*{Sperrvermerk}
			\thispagestyle{empty}
			Die Vorliegende Arbeit beinhaltet vertrauliche Informationen des Unternehmens \CompanyName. 
			Die Weitergabe oder Verwendung der Arbeit oder einzelner Informationen daraus ist in jeglicher
			Form untersagt. Dieses Verbot gilt nicht für Informationen, die weder aus dem Unternehmen \CompanyName\ stammen noch diese betreffen. Sonstige Ausnahmen bedürfen der schriftlichen Genehmigung des Unternehmens \CompanyName.
		\medskip
	\fi
}

%% End of File
	\newpage

	% da wir den Titel haendisch erstellt haben entfaellt der folgende Befehl:
	%\maketitle

	%\clearpage

	%\pagestyle{headings} 
	\pagestyle{fancy}   % Seite mit Kopf- und Fusszeilen darstellen
	\pagenumbering{roman}  % Umschalten auf Seitenzahlen in römischer Darstellung
	%\fancyfoot[R]{i}  % Seitennummer roemisch 1, hier per Hand eingetragen, weil es anders nicht funktioniert hat - hat sich erledigt


	\tableofcontents
	\newpage
	\listoffigures
	\newpage
	\listoftables
	
	\ifdefined\ThesisLanguageIsEnglish
		\printglossary[title=Glossary, toctitle=Glossary]
	\else
	%		\printglossary[title=Glossar, toctitle=Glossar]
		\printglossaries
	\fi
	
	% Allgemeine Akronyme
\chapter*{Abkürzungsverzeichnis}
\begin{acronym}[LONG]
	\acro{DUT}{Device under Test}
\end{acronym}
%% End of File %%
	
	\clearpage
	\pagenumbering{arabic}  % Umschalten auf Seitenzahlen in arabischer Darstellung
	%%%%% Einbinden der Dateien für die einzelnen Sektionen:
	% Verwendet man hier \include statt \input, beginnen neue Sektionen immer auf einer neuen Seite
	%\input{TeXFiles/Abstract}

	\fancyfoot[R]{\thepage}  % Seitennummer


	%\myName
	\chapter*{Abstract}
\label{ch:Abstract}




\newpage

\chapter*{Kurzfassung}
\label{ch:Kurzfassung}










%% End of File
	%\input{TeXFiles/004_LaTeXHowTo}      %%%%%% <<<<<< nach dem ersten Ansehen, bitte diese Zeile auskommentieren
	\chapter{Einführung}
\label{ch:Einfuehrung}
\lipsum[1-3]





\section{Zielstellung}
\label{sec:Zielstellung}
\lipsum[1-2]





\section{Aufbau der Arbeit}
\label{sec:Aufbau der Arbeit}
\lipsum[1]





%% End of File %%
	\chapter{Stand der Technik}
\label{ch:StandTechnik}





\section{Sektion}
\lipsum[1]




\subsection{Subsektion}
\lipsum[1]



\subsubsection{Subsubsektion}
\lipsum[1-5]






















%% End of File %%
	\chapter{Entwickelte Methode: Supereffiziente Methode zur Erstellung von optimierten xyz}
\chaptermark{Entwickelte Methode}  % short form for headlines on pages
\label{ch:EntwickelteMethode}





\section{Sektion}
\sectionmark{Sek} % short form for headlines on pages
\lipsum[1]




\subsection{Subsektion}
\lipsum[1]



\subsubsection{Subsubsektion}
\lipsum[1-5]







%% End of File %% 
	\chapter{Evaluation der Methode}
\label{ch:EvaluationDerMethode}





\section{Sektion}
\sectionmark{Sek} % short form for headlines on pages
\lipsum[1]




\subsection{Subsektion}
\lipsum[1]



\subsubsection{Subsubsektion}
\lipsum[1-5]
	\chapter{Ergebnisse}
\label{ch:Ergebnisse}





\section{Sektion}
\sectionmark{Sek} % short form for headlines on pages
\lipsum[1]




\subsection{Subsektion}
\lipsum[1]



\subsubsection{Subsubsektion}
\lipsum[1-5]
	\chapter{Résumé und Ausblick}
\label{ch:Zusammenfassung}

\lipsum[1-3]


% End of File %%

	\begin{appendix}
		\chapter{Anhang}
\label{ch:Anhang}






%% End of File %%
	\end{appendix}

	\clearpage

	
	%\printglossary[type=acronym,title=Abbreviations]
	%\addcontentsline{toc}{section}{Abbreviations} 
	%\addcontentsline{toc}{chapter*}{Verzeichnis: Akronyme} 
	\clearpage

	%Falls man die Überschrift des Literaturverzeichnisses aendern moechte, geht das durch Verwendung der folgenden Zeile:
	\renewcommand*{\refname}{Literaturverzeichnis}

	%\bibliographystyle{plain} % Nummern in eckigen Klammern
	%%\bibliographystyle{alpha} % Anfangsbuchstaben Erstautor und Jahr

	% Wenn man folgenden Stil verwenden will, dann muss
	% \usepackage{harvard} vor \begin{document} aktiviert werden:
	%\bibliographystyle{agsm} % (Autoren in runden Klammern)
	%\newpage
	\clearpage
	\phantomsection % Da sonst falsche Verlinkung im Inhaltsverzeichnis
	\printbibliography[heading=bibintoc]    % Ohne Kapitelnummer /-buchstabe
	
	%\printbibliography[heading=bibnumbered]    % Mit Kapitelnummer /-buchstabe
	
	%*******************************************************
% Declaration
%*******************************************************

\ifdefined\ThesisLanguageIsEnglish
	\chapter*{Declaration}
		\thispagestyle{empty}
		Hereby I assure that I have written the presented thesis independently and without any third party help and that I have not used any other tools or resources than the ones mentioned/referred to in the thesis.
	\medskip

	\noindent
		Any parts of the thesis that have been taken from other published or yet unpublished works in terms of their wording or meaning are thoroughly marked with an indication of the source. 
	\medskip

	\noindent
		All figures in this thesis have been drawn by myself or are clearly attributed with a reference. 
	\medskip

	\noindent
		This work has not been published or presented to any other examination authority before.
		I am aware of the importance of the affidavit and the consequences under examination
		law as well as the criminal consequences of an incorrect or incomplete affidavit.
	\medskip


\else
	\chapter*{Eidesstattliche Erklärung}
		\thispagestyle{empty}
			Ich versichere hiermit, dass ich die vorliegende Arbeit selbständig verfasst und keine anderen als die im Literaturverzeichnis angegebenen Quellen benutzt habe.
		\medskip

		\noindent
			Alle Stellen, die wörtlich oder sinngemäß aus veröffentlichten oder noch nicht veröffentlichten Quellen entnommen sind, sind als solche kenntlich gemacht.
		\medskip

		\noindent
			Die Zeichnungen oder Abbildungen in dieser Arbeit sind von mir selbst erstellt worden oder mit einem entsprechenden Quellennachweis versehen.
		\medskip

		\noindent
			Diese Arbeit ist in gleicher oder ähnlicher Form noch bei keiner anderen Prüfungsbehörde eingereicht worden. 
		\bigskip
\fi

Musterhausen, \ThesisDeliveryDate

\smallskip
%\includegraphics[width=5cm]{Figures/Unterschrift/YourSignature_w.png}
\hfill {\raisebox{-2ex}{\ \ }} \\ 
\vspace*{-1.7cm}
\begin{flushright}
	\begin{tabular}{m{5cm}}
		\\ \hline
		\centering \textbf{\myName}\\
	\end{tabular}
\end{flushright}

%% End of File
\end{document}

%% End of File